\documentclass[12pt]{article}
\usepackage[T1]{fontenc}
\usepackage[french]{babel}
\usepackage{datetime}
\usepackage{subcaption}
\usepackage[letterpaper, margin=1.9cm]{geometry}
\usepackage{graphicx}
\usepackage{epstopdf}
\usepackage{siunitx}
\graphicspath{ {../figures} }

\title{MEC6616 --- LAP1}


\author{Aurélie Grebe\\Noémie Herbinière}
\newdate{duedate}{1}{09}{2026}
\date{\displaydate{duedate}}

\begin{document}
\maketitle

\section{Exemple 4.1}

L'exemple 4.1 de Versteeg et al. est un problème stationaire de transfert de
chaleur d'un bar isolé dont la température sur d'un bord est fixé
à \SI{100}{\celsius} et l'autre à \SI{500}{\celsius}. Il s'agit donc des conditions
aux limites de Dirichlet.


\begin{figure}
	\centering
	\begin{subfigure}[t]{0.45\textwidth}
		\includegraphics[width=\textwidth]{ex4_1.png}
		\caption{Résultats analytique et par volume finis de l'exemple 4.1}
		\label{fig:ex4_1}
	\end{subfigure}
	\hfill
	\begin{subfigure}[t]{0.45\textwidth}
		\includegraphics[width=\textwidth]{epsilon4_1.png}
		\caption{Résultats analytique et par volume finis de l'exemple 4.1}
		\label{fig:ex4_1}
	\end{subfigure}
\end{figure}
\section{Exemple 4.2}
\begin{figure}
	\centering
	\begin{subfigure}[t]{0.45\textwidth}
		\includegraphics[width=\textwidth]{ex4_2.png}
		\caption{Résultats analytique et par volume finis de l'exemple 4.1}
	\end{subfigure}
	\hfill
	\begin{subfigure}[t]{0.45\textwidth}
		\includegraphics[width=\textwidth]{epsilon4_2.png}
		\caption{Résultats analytique et par volume finis de l'exemple 4.1}
	\end{subfigure}
\end{figure}
\section{Exemple 4.3}
\begin{figure}
	\centering
	\begin{subfigure}[t]{0.45\textwidth}
		\includegraphics[width=\textwidth]{ex4_3.png}
		\caption{Résultats analytique et par volume finis de l'exemple 4.1}
	\end{subfigure}
	\hfill
	\begin{subfigure}[t]{0.45\textwidth}
		\includegraphics[width=\textwidth]{epsilon4_3.png}
		\caption{Résultats analytique et par volume finis de l'exemple 4.1}
	\end{subfigure}
\end{figure}

	
\end{document}
